% Appendix — Key Code Listings

\section{ElasticNet Bootstrap Stability Selection (Python)}

The following listing implements the two-stage feature selection pipeline used in Experiments 4 and 6b. The first stage applies a hard mutual information threshold to remove features with negligible marginal association. The second stage applies ElasticNet bootstrap stability selection, retaining only features that are selected in more than 30 percent of bootstrap iterations.

\begin{verbatim}
from sklearn.linear_model import LogisticRegression
from sklearn.feature_selection import mutual_info_classif
from sklearn.utils import resample
import numpy as np

def mi_prefilter(X_train, y_train, threshold=0.002):
    """Stage 1: Remove features with MI below hard threshold."""
    mi_scores = mutual_info_classif(X_train, y_train,
                                    random_state=42)
    mask = mi_scores >= threshold
    return mask, mi_scores

def elasticnet_stability_selection(X_train, y_train,
                                   n_bootstraps=30, C=0.1,
                                   l1_ratio=0.5, threshold=0.30):
    """Stage 2: ElasticNet bootstrap stability selection."""
    n_features = X_train.shape[1]
    selection_counts = np.zeros(n_features)
    for _ in range(n_bootstraps):
        X_b, y_b = resample(X_train, y_train,
                            stratify=y_train, random_state=None)
        model = LogisticRegression(
            penalty='elasticnet', solver='saga', C=C,
            l1_ratio=l1_ratio, class_weight='balanced',
            max_iter=1000)
        model.fit(X_b, y_b)
        selected = np.abs(model.coef_[0]) > 0
        selection_counts += selected
    stability_scores = selection_counts / n_bootstraps
    stable_mask = stability_scores >= threshold
    # Verify p/n ratio is in the correct regime
    n_selected = stable_mask.sum()
    n_samples = len(y_train)
    assert n_selected / n_samples < 0.15, (
        f"p/n = {n_selected/n_samples:.2f} exceeds 0.15 limit")
    return stable_mask, stability_scores

# Two-stage pipeline usage
mi_mask, _ = mi_prefilter(X_train, y_train)
X_filtered = X_train[:, mi_mask]
stable_mask, scores = elasticnet_stability_selection(
    X_filtered, y_train)
X_selected = X_filtered[:, stable_mask]
\end{verbatim}

\section{Monte Carlo Dropout Uncertainty Quantification (PyTorch)}

The following listing implements the Monte Carlo Dropout predictor used in the explainability pipeline. Dropout layers are forced to training mode during inference to produce stochastic predictions. Predictive entropy is computed from the mean probability distribution across 30 forward passes.

\begin{verbatim}
import torch
import numpy as np

def mc_dropout_predict(model, x_temporal, x_spectral,
                       n_passes=30):
    """
    Approximate Bayesian posterior via MC Dropout.
    Returns mean probabilities, predictive entropy,
    and normalised confidence ratio.
    """
    # Force all dropout layers to training mode
    for m in model.modules():
        if isinstance(m, torch.nn.Dropout):
            m.train()
    probs = []
    with torch.no_grad():
        for _ in range(n_passes):
            logits, _, _ = model(x_temporal, x_spectral)
            prob = torch.softmax(logits, dim=-1)
            probs.append(prob.unsqueeze(0))
    probs = torch.cat(probs, dim=0)   # [N, B, T, C]
    mean_probs = probs.mean(dim=0)    # [B, T, C]
    eps = 1e-9
    entropy = -(mean_probs *
                torch.log(mean_probs + eps)).sum(-1)
    max_entropy = np.log(mean_probs.shape[-1])  # log(5)
    confidence_ratio = entropy / max_entropy
    return mean_probs, entropy, confidence_ratio

def confidence_tier(ratio):
    """Map normalised entropy ratio to three-tier label."""
    if ratio < 0.45:
        return "High"
    elif ratio < 0.72:
        return "Medium"
    else:
        return "Low"
\end{verbatim}

\section{Temporal-Spectral Fusion Model (PyTorch)}

The following listing presents the core components of the temporal-spectral fusion architecture: the Squeeze-and-Excitation block, the Adaptive Atrous Pyramid temporal encoder, and the top-level fusion model.

\begin{verbatim}
import torch
import torch.nn as nn

class SEBlock(nn.Module):
    """Squeeze-and-Excitation channel attention block."""
    def __init__(self, channels, reduction=16):
        super().__init__()
        self.pool = nn.AdaptiveAvgPool1d(1)
        self.fc = nn.Sequential(
            nn.Linear(channels, channels // reduction),
            nn.ReLU(),
            nn.Linear(channels // reduction, channels),
            nn.Sigmoid()
        )
    def forward(self, x):
        w = self.pool(x).squeeze(-1)
        w = self.fc(w).unsqueeze(-1)
        return x * w

class AdaptiveAtrousPyramid(nn.Module):
    """Multi-scale dilated temporal encoder."""
    def __init__(self, embed_dim=128):
        super().__init__()
        self.branches = nn.ModuleList([
            nn.Sequential(
                nn.Conv1d(1, embed_dim, kernel_size=3,
                          dilation=d, padding=d),
                nn.BatchNorm1d(embed_dim),
                nn.ReLU(),
                SEBlock(embed_dim)
            ) for d in [1, 2, 4, 8]
        ])
        self.gate = nn.Linear(
            len(self.branches) * embed_dim, embed_dim)
        self.pool = nn.AdaptiveAvgPool1d(1)

    def forward(self, x):
        # x: [B, 3000]
        x = x.unsqueeze(1)  # [B, 1, 3000]
        outs = [b(x) for b in self.branches]
        cat = torch.cat(outs, dim=1)  # [B, 4*D, T]
        cat_pooled = self.pool(cat).squeeze(-1)
        weights = torch.softmax(
            self.gate(cat_pooled), dim=-1)
        return (cat_pooled * weights)

class FusionSleepModel(nn.Module):
    """Temporal-spectral fusion transformer."""
    def __init__(self, embed_dim=128, n_heads=4,
                 n_layers=6, n_classes=5):
        super().__init__()
        self.temporal_enc = AdaptiveAtrousPyramid(embed_dim)
        self.spectral_enc = nn.Sequential(
            nn.Linear(34, embed_dim),
            nn.BatchNorm1d(embed_dim),
            nn.ReLU()
        )
        self.fusion = nn.Linear(embed_dim * 2, embed_dim)
        encoder_layer = nn.TransformerEncoderLayer(
            d_model=embed_dim, nhead=n_heads,
            dim_feedforward=512, dropout=0.2,
            batch_first=True)
        self.transformer = nn.TransformerEncoder(
            encoder_layer, num_layers=n_layers)
        self.head = nn.Linear(embed_dim, n_classes)

    def forward(self, x_temp, x_spec):
        # x_temp: [B, T, 3000], x_spec: [B, T, 34]
        B, T, _ = x_temp.shape
        h_t = self.temporal_enc(
            x_temp.view(B * T, -1)).view(B, T, -1)
        h_s = self.spectral_enc(
            x_spec.view(B * T, -1)).view(B, T, -1)
        h = self.fusion(torch.cat([h_t, h_s], dim=-1))
        h = self.transformer(h)
        return self.head(h), h_t, h_s
\end{verbatim}

\section{Explainability Output Feature Matrix}

Table~\ref{tab:platform} summarises the explainability outputs produced by the research pipeline, separated by intended audience.

\begin{table}[h]
\centering
\caption{Explainability output feature matrix by audience.}
\label{tab:platform}
\begin{tabular}{p{4.5cm}cc}
\toprule
Feature & Patient & Clinician \\
\midrule
Risk Score (0--100) & Yes & No \\
Feature Importance (\% influence) & Yes & No \\
Counterfactuals (plain language) & Yes & No \\
AI Advisor Chat & Yes & No \\
Personal Report Download & Yes & No \\
EEG Viewer (GradCAM regions) & No & Yes \\
Confidence Strip (per epoch) & No & Yes \\
Epoch Navigation (flagged) & No & Yes \\
Patient Directory & No & Yes \\
Upload (EDF / PT / CSV) & Yes & Yes \\
Sleep Study History & Yes & No \\
Personal Profile & Yes & No \\
\bottomrule
\end{tabular}
\end{table}

The patient interface strictly excludes entropy values, SHAP terminology, model accuracy, confidence scores, and EEG viewer access. The clinician interface provides full technical audit capability including GradCAM signal localisation, per-epoch confidence scoring, and flagged epoch navigation. This separation ensures that each user receives information appropriate to their role and technical background.
